% Experiment Template
% 1) Copy this file and rename it, e.g. experiments/experiment_04_firewall_rules.tex
% 2) In main.tex, add:
%    \includeexperiment{XX}{Experiment Title}{experiments/your_new_file_name}
% 3) Keep this file content as sections only (no \experiment command here)

\section{Theory}
Column transposition cipher arranges the plaintext into a table (grid) and then permutes the columns based on a key. For example, using \textbf{Jatiya Kabi Kazi Nazrul Islam University} in an $8 \times 5$ grid:

\begin{table}[H]
\centering
\begin{tabular}{|c|c|c|c|c|c|c|c|}
\hline
1 & 2 & 3 & 4 & 5 & 6 & 7 & 8 \\ \hline
J & a & t & i & y & a & \textvisiblespace & K \\ \hline
a & b & i & \textvisiblespace & K & a & z & i \\ \hline
\textvisiblespace & N & a & z & r & u & l & \textvisiblespace \\ \hline
I & s & l & a & m & \textvisiblespace & U & n \\ \hline
i & v & e & r & s & i & t & y \\ \hline
\end{tabular}
\caption{Characters arranged in an $8 \times 5$ table for Columnar Transposition}
\end{table}

Columns are then reordered according to a key (e.g. \textbf{3 1 4 2 5 6 8 7}), and the ciphertext is read row-by-row from the permuted table. Applying the inverse permutation reverts the ciphertext back to the original plaintext.

\section{Algorithm}
Encryption algorithm for Column transposition cipher
\begin{enumerate}
  \item Take the plaintext $T$ of length $n$ and a key array $K$ of length $k$.
  \item Initialize an empty string $E$ of length $n$.
  \item For each index $i$ from $0$ to $n-1$:
  \begin{enumerate}
    \item Calculate the source index in the plaintext as $j = K[i \pmod k] + \lfloor i / k \rfloor \times k$.
    \item Set the $i$-th character of the encrypted text: $E[i] = T[j]$.
  \end{enumerate}
  \item Return the encrypted string $E$.
\end{enumerate}

Decryption algorithm for Column transposition cipher
\begin{enumerate}
  \item Take the ciphertext $E$ of length $n$ and the same key array $K$ of length $k$.
  \item Initialize an empty string $D$ of length $n$.
  \item For each index $i$ from $0$ to $n-1$:
  \begin{enumerate}
    \item Calculate the destination index in the decrypted text as $j = K[i \pmod k] + \lfloor i / k \rfloor \times k$.
    \item Set the $j$-th character of the decrypted text: $D[j] = E[i]$.
  \end{enumerate}
  \item Return the decrypted string $D$.
\end{enumerate}

% \section{Flowchart (motherboard)}
% \begin{figure}[H]
%   \centering
%   \begin{autofitflowchart}
%     \node (start) [terminator] {Start};
%     \node (input) [io, below=of start] {Input};
%     \node (process1) [process, below=of input] {Process};
%     \node (decision1) [decision, below=of process1, yshift=-0.2cm] {Condition?};
%     \node (loopstep) [process, right=of decision1, xshift=1.2cm] {Update / Loop};
%     \node (output) [io, below=of decision1, yshift=-0.3cm] {Output};
%     \node (end) [terminator, below=of output] {End};

%     \draw [line] (start) -- (input);
%     \draw [line] (input) -- (process1);
%     \draw [line] (process1) -- (decision1);

%     \draw [line] (decision1.east) -- node[above]{No} (loopstep.west);
%     \draw [line] (loopstep.north) |- ([xshift=1.0cm]process1.east) |- (process1.east);

%     \draw [line] (decision1) -- node[right]{Yes} (output);
%     \draw [line] (output) -- (end);
%   \end{autofitflowchart}
%   \caption{Flowchart title}
% \end{figure}

\section{Code}
\begin{minted}{cpp}
#include<bits/stdc++.h>
using namespace std;

class column {
  vector<int> key;
public:
  column(int k);
  void showKey();
  string encrypt(string text);
  string decrypt(string text);
};
column::column(int k) {
  for (int i = 0; i < k; i++) {
    key.push_back(i);
  }
  for (int i = 0; i < k*10; i++) {
    int j = rand() % k;
    swap(key[i % k], key[j]);
  }
}

void column::showKey() {
  cout << "\nKey is: ";
  for (int i = 0; i < key.size(); i++) {
    cout << key[i] << " ";
  }
  cout << endl;
}

string column::encrypt(string text) {
  int k = key.size();
  int n = text.size();
  string encrypted(n, ' ');
  for (int i = 0; i < n; i++) {
    encrypted[i] = text[key[i % k] + (i / k) * k];
  }
  return encrypted;
}

string column::decrypt(string text) {
  int k = key.size();
  int n = text.size();
  string decrypted(n, ' ');
  for (int i = 0; i < n; i++) {
    decrypted[key[i % k] + (i / k) * k] = text[i];
  }
  return decrypted;
}


int main() {
  int k;
  cout << "Enter the key length: ";
  cin >> k;
  cin.ignore();
  column cipher(k);
  cipher.showKey();
  cout << "Enter Plaintext : ";
  string text;
  getline(cin, text);
  string encrypted = cipher.encrypt(text);
  cout << "Encrypted text: " << encrypted << endl;
  string decrypted = cipher.decrypt(encrypted);
  cout << "Decrypted text: " << decrypted << endl;
  return 0;
}
\end{minted}

\section{Result and Discussion - I/O}
Column transposition key is randomly generated in our program.
\begin{minted}{text}
Enter the key length: 5

Key is: 2 4 0 1 3 
Enter Plaintext : This is encrypted with some looping of the text.
Encrypted text: i Ths niseytcrp iedw oths omeligopnft o  ehet.x
Decrypted text: This is encrypted with some looping of the tex .
\end{minted}

\section{Conclusion}
We successfully implemented a column transposition cipher with random key generation. Since it is vulnerable to statistical attacks, it is typically used as one layer within more complex encryption systems.

